\documentclass[11pt,a4paper]{article}
\usepackage[utf8]{inputenc}
\usepackage{amsmath,amsfonts,amssymb}
\usepackage{geometry}
\geometry{margin=1in}
\usepackage{hyperref}
\usepackage[many]{tcolorbox}
\usepackage{enumitem}
\usepackage{fancyhdr}
\usepackage{tabularx}
\usepackage{booktabs}

\pagestyle{fancy}
\fancyhf{}
\fancyfoot[L]{Universitas Bengkulu - Teknik Elektro}
\fancyfoot[R]{\thepage}
\def\footrule{{\color{gray}\hrule}}

\title{\vspace{-2cm}\textbf{Rubrik Penilaian\\Mata Kuliah: Persamaan Diferensial}}
\author{Novalio Daratha \& Muhammad Arfan\\Program Studi Teknik Elektro\\Universitas Bengkulu}
\date{2026}

\begin{document}

\maketitle

\begin{tcolorbox}[colback=blue!5!white,colframe=blue!75!black,title=Deskripsi Umum]
Dokumen ini menguraikan rubrik dan komponen penilaian untuk mata kuliah Persamaan Diferensial. Penilaian mencakup Ujian Tengah Semester (UTS), Ujian Akhir Semester (UAS), serta \textit{Case Method} / \textit{Team-Based Project}.
\end{tcolorbox}

\section*{1. Ujian Tengah Semester (UTS) \& Ujian Akhir Semester (UAS)}
Penilaian untuk setiap soal pada UTS dan UAS didasarkan pada tiga komponen utama: pemodelan fisis, metode matematika/integrasi, dan hasil akhir. 

\begin{table}[h]
\centering
\renewcommand{\arraystretch}{1.5}
\begin{tabularx}{\textwidth}{|c|X|c|}
\hline
\textbf{Komponen} & \textbf{Kriteria Penilaian} & \textbf{Bobot} \\
\hline
Pemodelan Fisis & Pemahaman terhadap konsep dasar, pembuatan model fisika atau rangkaian listrik yang tepat dari masalah yang diberikan. & 20\% \\
\hline
Metode Matematika & Pemilihan metode penyelesaian yang tepat, langkah-langkah integrasi dan penurunan rumus yang sistematis. & 40\% \\
\hline
Hasil Akhir & Ketepatan hasil perhitungan akhir, termasuk satuan (jika ada), kondisi awal (IVP), dan interpretasi solusi. & 40\% \\
\hline
\end{tabularx}
\end{table}

\section*{2. Case Method / Team-Based Project}
Penilaian proyek berbasis tim berfokus pada analisis masalah nyata, kolaborasi anggota, dan kemampuan mempresentasikan hasil pemodelan diferensial untuk sebuah sistem.

\begin{table}[h]
\centering
\renewcommand{\arraystretch}{1.5}
\begin{tabularx}{\textwidth}{|c|X|c|}
\hline
\textbf{Komponen} & \textbf{Kriteria Penilaian} & \textbf{Bobot} \\
\hline
Analisis Masalah & Ketajaman analisis dalam merumuskan permasalahan riil menjadi persamaan diferensial secara logis. & 30\% \\
\hline
Solusi \& Simulasi & Penerapan metode analitik atau numerik yang benar, disertai dengan simulasi/visualisasi respons sistem. & 40\% \\
\hline
Kerja Sama Tim & Distribusi tugas yang merata, komunikasi yang baik, dan kontribusi nyata dari setiap anggota kelompok. & 15\% \\
\hline
Presentasi/Laporan & Kerapian laporan (seperti penggunaan \LaTeX), penyajian yang terstruktur, dan kemampuan mempertanggungjawabkan hasil. & 15\% \\
\hline
\end{tabularx}
\end{table}

\vfill
\begin{center}
\textbf{Dokumen ini bersifat acuan untuk memastikan objektivitas dan transparansi penilaian.}
\end{center}

\end{document}
